%!TEX encoding = UTF8
%!TEX root = 0-notes.tex

\chapter*{Préface}


\section*{Syllabus}

6 ECTS. 30h CM, 30h TD.

\subsection*{Description}

Cette UE vise à présenter les concepts élémentaires d'arithmétique et de dénombrement utiles pour le début de la licence en mathématiques.

\subsection*{Objectifs}

\begin{enumerate}
	\item
	Dénombrement élémentaire.
	\begin{enumerate}
		\item
		Cardinal d'un ensemble fini.
		Cardinal et opérations ensemblistes.
		Cardinal et applicatoins injectives, surjectives, bijectives.
		Cardinal d'un ensemble d'applications.
		Nombre de parties d'un ensemble.
		Fonction indicatrice.
		\item
		Introduction aux cardinaux infinis.
		Bijection entre ensembles.
		Dénombrabilité.
		Argument diagonal de Cantor.
		$X$ et $\mathcal{P}(X)$ n'ont pas le même cardinal.
		$\R$ est indénombrable.
		\item
		Arrangements, permutations, combinaisons (coefficients binomiaux), triangle de Pascal, formule du binôme.
		\item
		Formule du crible générale (application au dénombrement des dérangements, des surjections, etc.).
		\item
		Relation binaire sur un ensemble.
		Relation d'équivalence, partition en classes d'équivalence, quotient d'un ensemble par une relation d'équivalence (exemples sur des ensembles déjà connus).
		Relation d'ordre, partiel, total, exemples.
		\item
		Applications à des exemples de probabilités élémentaires finies (nombre de car favorables/nombre de cas total).
	\end{enumerate}
	\item
	Arithmétique élémentaire dans $\Z$.
	\begin{enumerate}
		\item
		Nombres entiers, écriture dans une base.
		\item
		Divisibilité, nombres premiers (infinitude, algorithme du crible).
		Division euclidienne (algorithme d'Euclide).
		\item
		PGCD et PPCM.
		Théorème de Bézout (et algorithme d'Euclide étendu), nombres premiers entre eux, lemme d'Euclide, lemme de Gauss.
		Équations diophantiennes $ax+by=c$.
		Décomposition en produit de nombres premiers.
		Application : pour $n\in\N$, $\sqrt{n}$ est soit un entier, soit irrationnel.
		\item
		Arithmétique modulaire (congruences).
		Petit théorème de Fermat.
		Théorème chinois des restes.
		\item
		Étude de $\Z/n\Z$ vu comme un anneau.
		Inversibles, $\Z/n\Z$ est un corps si et seulement si $n$ est premier.
		Réinterprétation du théorème de Bézout.
		Réinterprétation du petit théorème de Fermat (définition de l'indicatrice d'Euler, théorème d'Euler).
		Réinterprétation du théorème chinois des restes.
		\item
		Illustration par la cryptographie.
	\end{enumerate}
\end{enumerate}

\subsection*{Pré-requis nécessaires}

Pogramme de mathématiques de S1 (principalement Logique, raisonnements, ensembles et applications) et programme de mathématiques du lycée (a minima spécialité mathématiques de première).

\subsection*{Pré-requis recommandés}

Pogramme de mathématiques de S1 (principalement Logique, raisonnements, ensembles et applications) et programme de mathématiques du lycée (idéalement spécialité mathématiques de terminale, voire option mathématiques expertes.)

\section*{Structure du cours}

Construction d'un cours de mathématiques : Exemple -- Remarque -- Définition -- Lemme -- Proposition -- Théorème -- Démonstration -- Corollaire.

Voir \emph{Mechanica, volume 1}, Leonhard Euler, 1736 (\href{https://scholarlycommons.pacific.edu/euler-works/15/}{E15}), pour un exemple de cours de mathématiques en latin.

Wolframalpha permet de faire du calcul exact.

\begin{definition*}
Associe un mot de vocabulaire à un objet mathématique.
Un mot de vocabulaire peut avoir un sens différent en mathématiques qu'en langage vernaculaire.
C'est par exemple le cas de « ensemble », « couple », « rationnel », « réel », ...
\end{definition*}

\begin{lemme*}
	Résultat préliminaire, parfois technique, utile à la démonstration d'une proposition plus importante.
\end{lemme*}

\begin{proposition*}
	Énoncé vrai et important.
\end{proposition*}

\begin{theorem*}
	Énoncé vrai et fondamental.
\end{theorem*}

\begin{corollaire*}
	Résultat particulier, découlant souvent quasi directement d'un théorème ou d'une proposition.
\end{corollaire*}

%\newcounter{preface}
\begin{Exercise}[label=1]%[counter=preface]
	Exemple d'exercice d'application directe.
\end{Exercise}
\begin{Exercise}[difficulty=1, label=2]%, counter=preface]
	Exemple de problème.
\end{Exercise}
\begin{Exercise}[difficulty=2, label=3]%, counter=preface]
	Exemple de problème d'approfondissement.
\end{Exercise}
\shipoutExercise
\begin{Answer}[ref=1]
	Solution 1.
\end{Answer}
\begin{Answer}[ref=2]
	Solution 2.
\end{Answer}
\begin{Answer}[ref=3]
	Solution 3.
\end{Answer}
\begin{multicols}{3}
\shipoutAnswer
\end{multicols}
